In platform engineering, \enquote{capabilities} are often the center of discussion around a platform.
A capability here is a way of describing a service a platform offers.
A capability can be as simple as \enquote{runs a container} and as complex as \enquote{provisions a clustered database across regions, with constant \gls{wal} archiving to achieve \gls{pitr}}.
In the previous chapters, we learned about promises, which map very closely to what a capability is. \parencite{platformengineering-architects, cncfPlatformEngineering}
\par
The platform offers a service to a user, which can be directly mapped to a promise \parencite{promise-theory}:

\[ P \xrightarrow{+b} U \]

Where $P$ is the platform and $U$ indicates any user out of the group of users using the platform. This means our capability would be $b$ in this case.

\subsection{Promising impositions in a platform}
An agent making an imposition inherently cannot keep its own promise, only the agent on which it imposes behavior can. A system like Ansible is inherently imposition based, which means we need to find a way to map an imposition into a promise.
The simplest and most common way to achieve that is by making a promise to impose.

First we make Ansible $A$ impose a configuration $c$ on any Server $S_?$ contained in $S$.

\[ \pi _1: A \impose{c}{5} S_?(S) \]

Following that, Ansible can promise its own behavior:

\[ \sigma _1: A \xrightarrow{+\pi _1} P \]

Which the platform can then promise to consume:

\[ \sigma _2: P \xrightarrow{-\pi _1} A \]

This means Ansible makes a promise to the platform $P$ to impose $c$ on any Server $S_?$ in the list of servers $S$. The platform promises to use the promise, which allows the platform to depend on a promise made by Ansible.
We now have a promise, however, it is simply Ansible promising its own behavior to our platform.
The issue here is that our platform cannot promise Ansible's behavior, since this would be in violation of the principle of autonomy. \parencite{promise-theory}

Instead, we may promise an outcome and make it subject to another promise:

\[ \rho _1: P \xrightarrow{+D|\pi _1} U \]

This means our platform $P$ now promises to provide a desired state $D$ to our User $U$, subject to the behavior $\pi _1$ being fulfilled, which it can rely on due to $\sigma _1$ and $\sigma _2$.
Our capability in this case is now $D$, which means our platform now has the capability to use a promise Ansible has made to it to apply a configuration to an arbitrary server.
